\documentclass[11pt]{article}
\usepackage[margin=2.0cm]{geometry}
\usepackage{booktabs}
\usepackage{amsmath}
\usepackage{amssymb}
\usepackage{enumitem}
\usepackage{array}
\usepackage{listings}
\usepackage{xcolor}
\usepackage[T1]{fontenc}
\usepackage{lmodern}

\setlength{\parindent}{0pt}
\setlength{\parskip}{3.6pt}
\newcommand{\code}[1]{\texttt{#1}}
\newcommand{\hr}{\vspace{2pt}\hrule\vspace{6pt}}
\renewcommand{\arraystretch}{1.12}

\definecolor{kw}{RGB}{0,90,160}
\definecolor{cmt}{RGB}{110,110,110}
\lstdefinestyle{vhdl}{
  language=VHDL, basicstyle=\ttfamily\footnotesize,
  keywordstyle=\color{kw}\bfseries, commentstyle=\color{cmt}\itshape,
  showstringspaces=false, columns=fullflexible, breaklines=true,
  frame=leftline, framesep=6pt, xleftmargin=8pt, aboveskip=6pt, belowskip=6pt
}
\lstset{style=vhdl}

\begin{document}

\begin{center}
{\Large\bfseries EE\,362 --- Computer Organization Laboratory}\\[2pt]
{\large\bfseries Lab 4 --- Guided Examples: Solution}\\[2pt]
{\itshape Ex 1 tri-state and \code{'Z'} \quad$\cdot$\quad Ex 2 sign vs zero extension \quad$\cdot$\quad Ex 3 the conversion chain}\\[4pt]
Fall 2026 \quad$\cdot$\quad Lab TA: Gandhi Torres
\end{center}
\hr

\begin{center}
\begin{tabular}{clll}
\toprule
& circuit & new idea & feeds the assignment as \\
\midrule
1 & \code{tristate\_buf} & \code{'Z'}: a signal with \textbf{no driver} & \code{TriBuffer} \\
2 & \code{extend\_demo} & widening depends on what bits \textbf{mean} & \code{ImmedXtnd} \\
3 & \code{convert\_demo} & \code{std\_logic\_vector} $\to$ number $\to$ integer & \code{ALUsftMUX}, Lab 5 \\
\bottomrule
\end{tabular}
\end{center}

\vspace{4pt}\hr

\section*{Example 1 --- The third state}

\subsection*{1.1 Not a Boolean function}

\begin{center}
\begin{tabular}{c|c|l}
\toprule
\# & \code{Sel} & $Y$ \\
\midrule
0 & 0 & \code{'Z'} \quad not driving \\
1 & 1 & $A$ \quad driving \\
\bottomrule
\end{tabular}
\end{center}

\begin{lstlisting}
Y <= A when Sel = '1' else 'Z';
\end{lstlisting}

Everything built so far obeyed Boolean algebra: every signal was \code{'0'} or
\code{'1'}. This one does not.

\begin{quote}
\textbf{\code{'Z'} is not a value. It is the absence of a value} --- ``I am not
driving this wire at all.'' The output range is $\{0,1,Z\}$, and Boolean algebra
has only $\{0,1\}$, so no AND/OR/NOT expression can describe this component.
\end{quote}

\subsection*{1.2 Why \code{std\_logic} is a \emph{resolved} type}

If two drivers touch one wire, something must decide what the wire reads.
\code{std\_logic} carries a \textbf{resolution function} that combines all
drivers:

\begin{center}
\begin{tabular}{c||ccc}
\toprule
& \code{0} & \code{1} & \code{Z} \\
\midrule
\code{0} & \code{0} & \textbf{\code{X}} & \code{0} \\
\code{1} & \textbf{\code{X}} & \code{1} & \code{1} \\
\code{Z} & \code{0} & \code{1} & \code{Z} \\
\bottomrule
\end{tabular}
\end{center}

\begin{itemize}[leftmargin=1.6em,itemsep=1pt]
\item two drivers \textbf{agreeing} $\rightarrow$ that value;
\item one driving, one at \code{'Z'} $\rightarrow$ the driven value ---
\textbf{this is how a bus works};
\item two drivers \textbf{disagreeing} $\rightarrow$ \code{'X'}.
\end{itemize}

\begin{quote}
\textbf{So \code{'X'} in your waveform never means the simulator is confused. It
means two things drove one wire and disagreed.} In Labs 9 and 12 that is the
fastest diagnosis you own.
\end{quote}

\subsection*{1.3 Two things students get wrong here}

\textbf{(a) Contention is resolved per bit.} Put \code{0xAAAAAAAA} and
\code{0xBBBBBBBB} on the same 32-bit bus and the result is \textbf{not} all
\code{'X'}. Those patterns agree on some bits and disagree on others, so only the
conflicting bits go \code{'X'}. A check written as ``expect all X'' fails against
a correct simulator.

\textbf{(b) A floating bus is not zero.} With every driver disabled the wire is
all \code{'Z'}, not \code{0x00000000}. Reading a float gives \code{'Z'}, and any
arithmetic on it produces \code{'X'}.

\vspace{4pt}\hr

\section*{Example 2 --- Widening is two different operations}

\begin{center}
\small
\begin{tabular}{l>{\raggedright\arraybackslash}p{5.6cm}>{\raggedright\arraybackslash}p{5.6cm}}
\toprule
& what it does & when it is right \\
\midrule
\textbf{zero extension} & pad the top with \code{'0'} & the bits are \textbf{unsigned}, or a pattern \\
\textbf{sign extension} & pad the top with copies of the top bit & the bits are a \textbf{two's complement number} \\
\bottomrule
\end{tabular}
\end{center}

With a 4-bit input \code{0x8}: zero-extended it is \code{0x08} $= 8$;
sign-extended it is \code{0xF8} $= -8$. \textbf{Same input, different answers, and
neither is ``more correct''} --- the instruction decides.

\begin{lstlisting}
Yzero <= (7 downto 4 => '0')   & I;    -- top bits are zeros
Ysign <= (7 downto 4 => I(3))  & I;    -- top bits are copies of the sign bit
\end{lstlisting}

\textbf{The aggregate \code{(others => I(3))} is the idiom to learn.} It means
``every remaining element gets a copy of bit 3'' and works at any width without
being rewritten --- far better than typing \code{I(3) \& I(3) \& I(3) \& I(3)}.

\subsection*{Why this decides whether your CPU works}

\code{ADDI} sign-extends: its immediate is a signed offset, and
\code{ADDI r1,r1,-1} must actually subtract one. \code{LUI} zero-extends in this
CPU: its immediate is a bit pattern, not a number. \textbf{Get it backwards and
the Lab 12 loop counter counts the wrong way}, while the ALU and register file
both behave perfectly.

\vspace{4pt}\hr

\section*{Example 3 --- The conversion chain}

\code{std\_logic\_vector} is a \textbf{bag of bits}. It has no numeric meaning, so
VHDL will not add two of them --- it does not know whether you meant signed or
unsigned. You have to say.

\begin{center}
\code{std\_logic\_vector} $\longrightarrow$ \code{unsigned} / \code{signed}
$\longrightarrow$ \code{integer}
\end{center}

\begin{lstlisting}
to_integer(unsigned(Slv))                  -- bits as unsigned, then an integer
to_integer(  signed(Slv))                  -- bits as two's complement
std_logic_vector(to_unsigned(N, W))        -- and back again, at width W
\end{lstlisting}

Read it inside out: \code{unsigned(Slv)} \emph{reinterprets} the bits;
\code{to\_integer()} then converts.

\subsection*{Why you cannot avoid it}

\textbf{You cannot index a vector with a vector.} \code{Data(Shamt)} is
meaningless --- array indices are integers. So a shift amount, a register number
or a memory address arriving as bits must be converted first:

\begin{lstlisting}
signal amount : natural range 0 to 7;
...
amount <= to_integer(unsigned(Shamt));           -- convert ONCE
Bit0   <= Data(amount);                          -- index needs an integer
Shift  <= std_logic_vector(shift_left(unsigned(Data), amount));
\end{lstlisting}

\textbf{Convert once into a named signal, then use it.} Burying the conversion
inside a larger expression is where the unreadable errors come from.

\subsection*{The library rule, and why it is absolute}

\begin{lstlisting}
use ieee.numeric_std.all;        -- yes, always
-- NEVER: std_logic_arith, std_logic_unsigned, std_logic_signed
\end{lstlisting}

Those three are not IEEE standard, they conflict with each other, and mixing them
produces errors that read as though the compiler has lost its mind --- typically
``no feasible entries for operator'' on a line that is obviously fine.

\vspace{4pt}\hr

\section*{Into the assignment}

\begin{center}
\begin{tabular}{ll}
\toprule
example & becomes \\
\midrule
Ex 1 & \code{TriBuffer}: \code{Dout <= Din when Sel = '1' else (others => 'Z');} \\
Ex 2 & \code{ImmedXtnd}: extend by \code{Xsel}, \textbf{then} shift by \code{Xlsft} \\
Ex 3 & \code{ALUsftMUX}, and every index you write from Lab 5 onward \\
\bottomrule
\end{tabular}
\end{center}

\code{ImmedXtnd} is built in \textbf{two visible steps} for a reason:

\begin{lstlisting}
extended <= (31 downto 16 => Immed16(15)) & Immed16 when Xsel = '1'
       else (31 downto 16 => '0')         & Immed16;

Xtnd32   <= extended(29 downto 0) & "00" when Xlsft = '1'
       else extended;
\end{lstlisting}

\textbf{Extend first, then shift.} Both branches of both assignments produce a
value, so there is no latch --- compare Lab 3, Example 3. And the RTL Viewer shows
the two stages separately, which makes the order visible rather than a claim.

\begin{center}
\begin{tabular}{c|ccc|l}
\toprule
\# & \code{Xsel} & \code{Xlsft} & \code{Immed16} & \code{Xtnd32} \\
\midrule
0 & 0 & 0 & \code{0x8000} & \code{0x00008000} \\
1 & 1 & 0 & \code{0x8000} & \code{0xFFFF8000} \\
2 & 1 & 1 & \code{0x0002} & \code{0x00000008} \\
3 & 1 & 1 & \code{0xFFFF} & \code{0xFFFFFFFC} \\
\bottomrule
\end{tabular}
\end{center}

Check row 3 as a number: $-1 \times 4 = -4 =$ \code{0xFFFFFFFC}. \textbf{Sign
extension followed by $\times 4$ still multiplies the value by four}, negatives
included --- that is the arithmetic sanity check for this component.

\vspace{6pt}
\hr
{\small \textbf{Forward.} \code{TriBuffer} is what lets the register file and the
memory share one data bus in Lab 12; \code{ImmedXtnd} is what makes \code{ADDI},
\code{LUI}, \code{LW}, \code{SW} and the branches agree about what their
immediates mean; and the conversion chain appears in every component that indexes
anything.}

\end{document}
